% ML Experiment Triage: main report.
% Build with `make report`. Every number and figure is generated from the
% database by scripts/gen_report_assets.py; nothing here is typed by hand.
\documentclass[11pt,a4paper]{report}

\usepackage[T1]{fontenc}
\usepackage[utf8]{inputenc}
\usepackage{lmodern}
\usepackage[margin=2.6cm]{geometry}
\usepackage{booktabs}
\usepackage{tabularx}
\usepackage{microtype}
\usepackage{amsmath}
\usepackage{amssymb}
\usepackage{pgfplots}
\usepackage{xcolor}
\usepackage{caption}
\usepackage[hidelinks]{hyperref}

\pgfplotsset{compat=1.18}
\usepgfplotslibrary{fillbetween}

% The validated reference categorical palette, in its documented slot order.
\definecolor{seriesA}{HTML}{2A78D6}
\definecolor{seriesB}{HTML}{EB6834}
\definecolor{seriesC}{HTML}{1BAF7A}
\definecolor{seriesD}{HTML}{EDA100}
\definecolor{seriesE}{HTML}{E87BA4}
\definecolor{seriesF}{HTML}{008300}
\definecolor{baselinegrey}{HTML}{52514E}
\definecolor{gridgrey}{HTML}{E1E0D9}
\definecolor{axisgrey}{HTML}{C3C2B7}

\captionsetup{font=small, labelfont=bf}
\setlength{\parskip}{0.4em}

\input{tables/facts.tex}

\title{\textbf{ML Experiment Triage}\\[0.4em]
\large Ranked, significance tested comparison of machine learning training runs}
\author{Olajide Badejo}
\date{5 August 2026}

\begin{document}
\maketitle

\begin{abstract}
\noindent
Comparing training runs by looking at curves fails in four compounding ways: the
final point is mostly noise, curve values are autocorrelated enough to break the
usual tests, run to run variance is routinely larger than the effects being
claimed, and several metrics are compared at once. This report describes a tool
that replaces that practice with a stated hypothesis, a permutation test whose
assumptions the data actually satisfies, a two gate regression rule with false
discovery control, and, most importantly, a calibration suite that measures the
tool's own error rates rather than assuming them.

The measured type I error is 4.53 percent against a nominal 5, with power of
96.25 percent on a large effect. The report also publishes what the weaker single
run mode costs: on synthetic runs with realistic seed variance and a true effect
of exactly zero, comparing one run against one run calls the difference
significant on up to 87 percent of cases, where the seed replicated mode holds at
5. That number, rather than any feature, is the argument for the design.
\end{abstract}

% The document wide paragraph skip is added between every contents line, which
% pushed the single appendix entry onto a second page on its own. Zeroing it
% inside this group keeps the contents on one page and leaves body text alone.
{\setlength{\parskip}{0pt}\tableofcontents}

% ---------------------------------------------------------------------------
\chapter{Introduction}

The routine this project replaces is familiar. Two runs, a TensorBoard tab, one
loss curve slightly below the other at the right hand edge, and a decision made
on it. It is quick, it feels evidential, and it is close to worthless.

The aim was to build something that answers a sharper question: which runs beat
the baseline, by how much, with what statistical confidence, and, crucially,
with the tool's own calibration proven by tests rather than asserted in a README.
The last clause is the one that shaped everything else. A tool that produces p
values is only useful if its p values mean what they say, and the only way to
know that is to run it thousands of times against data whose truth you set
yourself and count how often it is wrong.

\section{Objectives}

\begin{enumerate}
  \item Parse TensorBoard event files, CSV and JSONL logs into one experiment
    model, stored in SQLite so nothing is ever reparsed.
  \item Compare any run against a baseline with a permutation test on a final
    window statistic, reporting effect size, confidence interval and p value, in
    two clearly labelled modes.
  \item Flag regressions on two gates, statistical and practical, with
    Benjamini Hochberg correction across metrics and ranking by severity.
  \item Rank hyperparameter sensitivity, always reporting the sample size behind
    each correlation.
  \item Treat statistical calibration as a first class test.
  \item Produce a self contained HTML report and reproducible PDFs.
\end{enumerate}

\section{What this report contains}

Chapter~\ref{ch:background} sets out why curve comparison needs care.
Chapter~\ref{ch:methodology} defines the statistic and both modes precisely, and
states what each can and cannot claim. Chapter~\ref{ch:implementation} covers the
engineering. Chapter~\ref{ch:results} gives the measured calibration, the demo
verdict table and the sensitivity analysis. Chapter~\ref{ch:discussion} discusses
the limitations honestly, and Chapter~\ref{ch:conclusion} concludes.

% ---------------------------------------------------------------------------
\chapter{Background}
\label{ch:background}

\section{Autocorrelation}

Consecutive evaluations of a training run are correlated: the model at step 5000
is nearly the model at step 4999. A test that treats a window of curve points as
independent observations computes a standard error far too small and reports
significance that is not there. This is not a small inefficiency. It is the
difference between a test that controls its error rate and one that does not.

This is why the permutation test here uses the \emph{run} as its unit of
analysis, and why the single run mode has to permute blocks rather than points.

\section{Multiple comparisons}

Testing eight metrics at $\alpha = 0.05$ gives roughly a one in three chance of at
least one false positive when nothing has changed. Any tool that compares several
metrics at once and does not correct for it is manufacturing findings.

Benjamini and Hochberg's step up procedure controls the false discovery rate and
is the right instrument here, because training metrics are strongly correlated
and Bonferroni assumes worst case dependence. On the fifteen p values in the
original 1995 paper, Bonferroni rejects three where the step up procedure rejects
four; that difference in power is real, and this project's implementation is
tested against exactly that example.

\section{Seed variance}

Two identical configurations trained with different seeds land in different
places, and in deep learning that spread is routinely larger than the effect
sizes claimed from single runs. Bouthillier and colleagues set this out in detail
in their MLSys 2021 accounting of variance in benchmarks, and it is the finding
this tool is built around.

The consequence is sharper than "be careful". A comparison with one run per side
contains \emph{no information} about run to run variance. It cannot separate the
effect of a change from the effect of a seed, not imprecisely but in principle.
Section~\ref{sec:weakcost} measures what that costs.

% ---------------------------------------------------------------------------
\chapter{Methodology}
\label{ch:methodology}

\section{The statistic}

For a tagged series of $n$ values, smooth with a centred moving average of width
nine, then take the mean of the final window: the last
$k = \max(20, \lceil 0.1n \rceil)$ points. The full specification as used is:

\begin{quote}\small\factstatistic\end{quote}

The window mean has roughly $\sqrt{k/\tau}$ times less noise than a single final
point, where $\tau$ is the integrated autocorrelation time. Using the whole curve
instead would mix how fast a run got going with how good it got, answering
neither question.

\section{Mode one: seed replicated}

\textbf{Null hypothesis.} The condition label is exchangeable across runs; the
final window statistics of the two conditions are draws from one distribution.

\textbf{Test.} Pool the $m+n$ run level statistics, and over every split into
groups of size $m$ and $n$ compute the difference of group means. The p value is
the fraction at least as extreme in absolute value as the observed difference.
With five seeds a side there are $\binom{10}{5} = 252$ arrangements, enumerated
exhaustively, so the p value is exact rather than estimated.

\textbf{Why this is the strong claim.} Seed variance sits inside the null
distribution rather than being assumed away. If two conditions differ by less
than the runs of one condition differ among themselves, the permutation
distribution is wide and the p value is large, which is correct.

\textbf{The design limit, reported not hidden.} With three seeds a side there are
twenty distinct assignments, so the smallest attainable two sided p value is 0.1
and no result can clear $\alpha = 0.05$. The tool computes this bound, warns when
it exceeds alpha, and classifies such findings as inconclusive rather than as no
change. A verdict of no significant difference from a design that could not have
found one is not a finding.

\section{Mode two: single run window block}

\textbf{Null hypothesis.} The two final windows are blockwise exchangeable: both
are segments of the same stationary process.

\textbf{Test.} Split each final window into contiguous blocks of three times the
estimated integrated autocorrelation time, estimated on each window separately,
and permute the block labels between the runs.

\textbf{Precondition, enforced.} The window must hold at least eight such blocks.
Where it does not, the tool raises an error naming the remedy rather than
shortening the block to fit. Section~\ref{sec:calibration} explains why that
distinction is not fastidiousness.

\textbf{What it cannot claim.} Its null is that these two windows come from the
same process, which is not the question anyone is really asking. Every output
from this mode is labelled a weaker claim.

\section{Two gates}

A finding is a regression only when both hold:

\begin{quote}\small\factgates\end{quote}

Direction is inferred from the tag name and can be overridden. Findings are
ranked by severity, which scales the size of the harm by confidence and halves
the result for the weaker mode, so a weaker claim never outranks a stronger claim
of the same size.

\section{Sensitivity}

Spearman rank correlation between each numeric hyperparameter and the final
window statistic, with the run count printed beside every correlation. Rank
correlation rather than ANOVA, because sweeps are unbalanced and hyperparameter
effects are usually monotone but strongly nonlinear.

% ---------------------------------------------------------------------------
\chapter{Implementation}
\label{ch:implementation}

\section{Model and storage}

Every parser produces one \texttt{Experiment}: an identifier, a configuration,
and a set of tagged step series. Everything downstream reads that model and
nothing else, which is what lets a TensorBoard run and a JSONL log be compared
without the analysis knowing where the numbers came from. The unit suite asserts
this directly: four committed fixtures encoding one identical series in four
containers parse to numerically identical experiments.

Series are held and stored as float32. Holding float64 in memory while writing
float32 would make a store round trip lossy in a way no test could honestly call
exact; fixing the model at the storage precision makes the round trip a bit for
bit assertion.

Ingest fingerprints each source on path, size and modification time. An unchanged
source is skipped, a changed one replaces its rows in a single transaction, and
an interrupted ingest loses only the run in flight. Content hashing was rejected
because hashing a directory of event files costs about as much as parsing it.

\section{Demo sweep}

The demo generates \factruns{} runs across \factconditions{} conditions, in three
log formats at once, from fixed seeds. Table~\ref{tab:runs} lists them.

\begin{table}[htbp]
\centering
\small
\input{tables/runs.tex}
\caption{The synthetic demo sweep. Spread is the standard deviation across seeds
within a condition, and it is the quantity every claim in the next chapter has to
overcome.}
\label{tab:runs}
\end{table}

The database holds \factpoints{} points across \factseries{} series in
\factdatabasekb{}~KB, a compression ratio of \factcompressionratio{} on the
series blobs.

% ---------------------------------------------------------------------------
\chapter{Results}
\label{ch:results}

\section{Measured calibration}
\label{sec:calibration}

This is the section the rest of the report depends on.

\begin{table}[htbp]
\centering
\small
\input{tables/calibration.tex}
\caption{Measured error rates against the phase gates. Every case is driven from
a fixed seed, so these numbers are deterministic and reproduced by
\texttt{make test-stats} in about a minute.}
\label{tab:calibration}
\end{table}

Checking the error rate at a single threshold can be passed by a test that is
wrong in two compensating directions, so the suite also checks that the null
distribution of p values is uniform.

\begin{table}[htbp]
\centering
\small
\input{tables/uniformity.tex}
\caption{Rejection rates under the null at four thresholds. A calibrated test is
right everywhere, not only at 0.05.}
\label{tab:uniformity}
\end{table}

\subsection{What calibration testing changed}

The first full run of the suite measured the window block mode at a type I error
of \textbf{12 percent} against a nominal 5. Two causes, both found by measurement
rather than by reading the code.

First, the autocorrelation time was estimated on the two windows concatenated.
Joining two independent runs end to end puts a step change at the join, which the
estimator reads as long range dependence, inflating the estimate by roughly a
factor of two.

Second, and more consequentially, the block length was one autocorrelation time.
At one $\tau$ the block means are still visibly correlated with their neighbours.
The observed split is contiguous, one run's blocks against the other's, but a
permuted split scatters neighbouring blocks across both groups, breaking that
correlation and producing a null distribution narrower than the truth. A null
that is too narrow is exactly how a test becomes anticonservative.

The fix was to estimate $\tau$ per window and set the block at three times it,
chosen by measuring type I error across window lengths from 1200 to 8000 steps
and autocorrelation coefficients from 0 to 0.9. That created a third decision:
once the block is three $\tau$, a short run may not contain eight of them, and
shortening the block to fit is precisely the failure above. So the mode refuses.
A quietly miscalibrated p value is worse than no p value, because it spends
credibility that nothing later can recover.

\section{The cost of the weaker mode}
\label{sec:weakcost}

Table~\ref{tab:weakmode} is the most useful number this project produced. Both
modes are run on synthetic runs carrying realistic seed variance, with a true
effect of \emph{exactly zero}. Every rejection is a false positive.

\begin{table}[htbp]
\centering
\small
\input{tables/weak_mode.tex}
\caption{False positive rates when the true effect is zero and seed variance is
present. Comparing one run against one run, on data where nothing is different,
calls it a significant difference most of the time.}
\label{tab:weakmode}
\end{table}

This is not a defect in the implementation. It is what the question ``are these
two curves different'' can answer when it has no information about how much a
curve moves between seeds. The tool still offers the mode, because sometimes one
run is all there is, and labels every result from it accordingly.

\section{Demo verdicts}

\begin{table}[htbp]
\centering
\scriptsize
\input{tables/verdicts.tex}
\caption{All \factcomparisons{} comparisons against \factbaseline{}, ranked by
severity: \factregressions{} regressions and \factimprovements{} improvements.
Rows marked $\dagger$ come from the weaker single run mode. Permutation seed
\factpermutationseed{}.}
\label{tab:verdicts}
\end{table}

The tool identifies \factbestcandidate{} as the best condition, which is the
ground truth the sweep was built with, and does not flag the condition whose true
effect is real but below the practical threshold.

One row deserves a second look. The single seed condition reports an adjusted p
of 0.0006, far smaller than the five seed conditions with comparable effects,
which sit between 0.016 and 0.032. That is the anticonservatism of
Table~\ref{tab:weakmode} appearing in the demo itself rather than only in the
calibration suite.

\begin{figure}[htbp]
\centering
\input{figures/val-loss-curves.tex}
\caption{Validation loss. Each line is the mean across seeds of a smoothed curve;
the band is the full range across seeds of that condition. Where two bands
overlap, the conditions are not separated by more than their own run to run
noise, whatever the lines appear to do.}
\label{fig:loss}
\end{figure}

\begin{figure}[htbp]
\centering
\input{figures/val-accuracy-curves.tex}
\caption{Validation accuracy for the same runs.}
\label{fig:accuracy}
\end{figure}

\section{Sensitivity}

\begin{table}[htbp]
\centering
\small
\input{tables/sensitivity.tex}
\caption{Spearman rank correlation of each swept hyperparameter against each
metric, over all \factruns{} runs.}
\label{tab:sensitivity}
\end{table}

Table~\ref{tab:sensitivity} demonstrates a real limitation rather than hiding it.
Learning rate drives by far the largest effects in this sweep, and shows the
\emph{weaker} correlation. The reason is that the sweep gives learning rate an
optimum in the middle of its range: 0.0003 is worse than 0.001, 0.003 is the best
of the four and 0.01 is by a distance the worst. That is a U shape, and the rank
correlation of a U shape is near zero however large the effect. Batch size, swept
monotonically over the same runs, correlates more strongly while driving a
smaller effect.

A near zero rank correlation here means \emph{not monotone}, never \emph{no
effect}. The tool states this wherever the table appears.

% ---------------------------------------------------------------------------
\chapter{Discussion}
\label{ch:discussion}

\section{What this tool is honest about}

Three limitations are documented, tested and printed rather than smoothed over:
the single run mode cannot see seed variance and says so; rank correlation is
blind to non monotone effects and the demo demonstrates it; and sensitivity is
association across a sweep, not a controlled effect.

Documenting a limitation is cheap. Building the demonstration of it into the
demo, and pinning it with a test named for it so nobody later ``fixes'' the
inconvenient number, is the part that took effort and is the part worth having.

\section{What I would do differently}

The Welch interval reported alongside the seed replicated p value leans on a
normal approximation the p value itself does not need. It is named in the output
for that reason. A percentile bootstrap over five seeds is known to be poor, so
the alternative was not obviously better, but a properly studentised bootstrap
would be more consistent.

The severity ranking is a defensible heuristic rather than a derived quantity.
Halving the weaker mode encodes a judgement about how much less a weak claim is
worth, and that number is not measured.

\section{Future work}

A nonlinear dependence measure alongside Spearman, such as distance correlation,
would see the U shape that rank correlation misses. A sequential mode that
decides how many further seeds are needed to resolve an inconclusive comparison
would be more useful than any of the above, since the most common honest verdict
this tool gives is that there is not enough data yet.

% ---------------------------------------------------------------------------
\chapter{Conclusion}
\label{ch:conclusion}

The tool does what it set out to do: it takes a directory of runs in three log
formats, and returns a ranked table saying which beat the baseline, by how much,
with what confidence, and under which of two clearly separated claims.

What makes it worth trusting is not any of that. It is that its error rates were
measured rather than assumed, that the measurement found a real defect running at
better than twice the nominal error rate, and that the fix was to make the tool
refuse rather than to widen the gate. The calibration suite is the product; the
comparison is what it certifies.

\appendix
\chapter{Reproducing this report}

\begin{verbatim}
make env      # Python 3.13 venv, pinned dependencies
make test     # full suite including the calibration gates
make demo     # regenerate the sweep, database and HTML report
make all      # everything, including this PDF
\end{verbatim}

Every table and figure above was generated by
\texttt{scripts/gen\_report\_assets.py} from the database, at permutation seed
\factpermutationseed{}. No number in this document was typed by hand.

\begin{thebibliography}{9}
\bibitem{bh1995} Y. Benjamini and Y. Hochberg. Controlling the false discovery
  rate: a practical and powerful approach to multiple testing. \emph{Journal of
  the Royal Statistical Society, Series B}, 57(1):289\,to\,300, 1995.
\bibitem{bouthillier2021} X. Bouthillier et al. Accounting for variance in
  machine learning benchmarks. \emph{Proceedings of Machine Learning and Systems
  (MLSys)}, 2021.
\bibitem{good2005} P. Good. \emph{Permutation, Parametric, and Bootstrap Tests of
  Hypotheses}, 3rd edition. Springer, 2005.
\bibitem{phipson2010} B. Phipson and G. K. Smyth. Permutation p values should
  never be zero. \emph{Statistical Applications in Genetics and Molecular
  Biology}, 9(1), 2010.
\end{thebibliography}

\end{document}
